\chapter{Dialog}

\section{Introduction and Scope}



\section{Politics}

\paragraph{AI is a Net Positive*}

\section{Philosophy}

\paragraph{Who Should I Become(part 1, Broad)}

Q: Before we continue, let us narrow the question down to a pieces. Come up with various categories and sub categories.

R: ok

\begin{itemize}
  \item Presence
  \begin{itemize}
    \item Physical Presentation
    \item Facial Expressions
    \item Tone of Voice
    \item Body Language
    \item Tone of ``warmth''
    \item Humor and Levity
    \item Professionalism
  \end{itemize}
  \item Speech
  \begin{itemize}
    \item Vocabulary and Precision
    \item Technical Explanation
    \item Academic Writing Structure
    \item Cross-audience communication
    \item Socratic Dialogue
  \end{itemize}
  \item Conversation and Social Interactions
  \begin{itemize}
    \item Core skills
    \item Networking
    \item Situation scripts
  \end{itemize}
  \item Mind and Cognition
  \begin{itemize}
    \item Cognative Habits
    \item mathematical and physical thinking
    \item Problem solving
    \item Desicion making
    \item Sustained and Consistent focus
    \item Observation
    \item Bayes thinking
    \item Emotional Regulation
  \end{itemize}
  \item Character...
\end{itemize}

Q: Ok, that is enough. We don't need more than that. What would you say is the most fundemental on that list?

R: Character, our morals dictate our actions, how we interprate actions, and much more. It influences everything else while being independed of everything else.

Q: Good, we have already discussed that in detail in this book, what would you say the next most fundemental? For the purposes of this exersices we will use fundemental to mean building blocks towards our conversation, not any other meaning you could describe.

R: For that, I would have to say that presence. As your goal of presence effects sets your goals for speech and conversation. Mind is largely independent, but it is also independent of the others. More orthogonal than fundemental.

Q: I accept your argument. To further constrain this conversation, let us define what you mean by presence.

R: I mean who I wish to be processed as. In a general sense.

Q: Does that fit with your above categories

R: Sadly not, but it can be useful in other situations. I would go back an reorder our conversation, but let us move ahead.

Q: Wouldn't that challenge the exercise?

R: No, the point of the beginning was simply to define what we should talk about, now we know. Furthering it could be of use, but we can do it later without such hidden truths coming to trouble in this conversation.

Q: Ah, so, now define our conversation.

R: While I find the whole purpose is to change internally for the sake of overcoming myself. Something I have over looked is others perception of me. Thus I should examine that. Here we will define broadly who I want to be perceived to be.

Q: Why?

R: Why I haven't thought of it before or why it matters?

Q: The latter.

R: Very simply, modern society is... well a society. To succeed and achieve my broader goals I rely on others to see me, to help, to teach, and to let me lead when the time comes. It is also as I have mentioned a human's duty to be pro-social within their society. Having a defining oneself so others can see allows other to allow them to do those things much more easily.

Q: I could certainly examine that short statement further, but let us move one. for an experiment, let us start with writing some general ideas who you want to be seen as.

R: Rigorous, consiouncious(in all of its definitions), conscious, questions assumptions, actively helpful, kind, helpful, friendly, obedient, reverent, independent but still a leader, curious, broad minded, disciplined, consistent, deep, ``always'', precise, high standards, active, confident, purposefull, makes others more capable, composed, long-term vision, and principled.

Q: Alright, we will reconvine in the morning.

R: Alright, I have returned.

Q: Looking at the above, in broad strokes what would you say the image you want to convey to be?

R: Competent, in control, and helpful

Q: Do any constraints arise?

R: Being seen as competent, in control, and all of the related properties tends to lean towards stoic and general neutral line. While being helpful and a friendly would obviously lean towards being warm and ``friendly.''

Q:Remember, whenever paradoxes arise, test you assumptions, as contradictions do not exist.

R: Yes, there is no true paradox, only the question of picking a point on the spectrum and knowing when to move about that spectrum.

Q: Why is this?

\paragraph{My use of AI}

Q: to question the use of AI, let us first constrain ourselves to
\section{Physics}


\paragraph{The Laws of Physics are the same everywhere}
Q: To examine the claim, we must first define out terms. So, if you would please explain what you mean by the Laws of Physics and ''are the same.''

R: Well, that is quite simple. The Laws of physics are the underlying mechanisms and devices that govern the Universe. Whether that be symmetries, quantum fields, structure, or something else has little barring on this topic. As for are the same, simply that there exists a spacial invariants across the universe.

Q: Why does the nature of the laws of physics have little baring on the conversation at hand?

R: All of the listed above have the same observable nature by definition. It is not like how the electron behaves differently than the atom, but rather more like the quantum field fluctuations and electron. If they had easily distinguishable differences that can be examined scientifically, such discussions would not be had.

Q: Ok, I accept your premises, now on to the claim itself. How exactly would one be able to verify that the laws of physics are spatially invariant if we use our current programme to continue research into the new?

R: While I consider that it is quite difficult, and while my axiom should still be held, a general skepticism is perfered. Now how to look out for it, modern advancements allow us to see areas separated by great distances. Inconsistencies accross different regions of space would be quite easy to observe. Allowing one to question held assumptions.

Q: What about areas beyond our observation?

R: No objective way to prove it, but science is not about unprovable ``what-ifs,'' it examines observables. As mentioned, it is taken as an axiom unless falsified.

Q: Yes, until falsified. You make it sound just like another theory. Yet, is it not common that another criteria for a theory would be to adhere to this axiom?

R: Yes, just as any theory must adhere to modern scientific consencous, until reason to doubt such underlying theories.

Q: I realize something as we talk. When referring to spacial changes in physics, you take a more conservative view than I. Focusing on the idea of variable pockets of contiuum with slight changes. Alas, there could be massive changes. Changes in dimensions, symmetries, or even fundamental properties like momentum transfer.

R: As for fundemental properties, they are simply consequences of more fundemental properties like dimensions and symmetries. Possible changes have been thoroughly examined so one would be able to observe them and theoretically figure it out.

Q: What if the changes are to logic and mathematics? Rather than symmetry.

R: That is absured.

Q: Why is that?

R: No possible mechanism could possibly do that, as logic cannot be effected by mechanisms.

Q: How could you know that? Do you have omniscient knowledge?

R: No, but deductive reasoning can show that math is divorced from physical reality and thus cannot be effected by physical mechanisms. The idea of logic being capable of changing by external forces is paradoxical. If logic changes, that would simply be a consequence of other logic, thus consistent with logic. Only new.

Q: So there could be new logic that we have no knowledge of?

R: Yes, but, such idyllic daydreaming is unrealistic.

Q: Why, what makes it any less unlikely than anything else?

R: There is no possible evidence of it.

Q: Would there be evidence of it?

R: Yes, changes in the laws of physics over time or space.

Q: You assume you understand this unknown, how could you? How could you know you could identify it with any for of accuracy? What makes you think the current mysteries aren't that our current understanding it reliant of completely false axioms?

R: Such avenues are being explored.

Q: In very limited forms. Not in the forms I have been suggesting.

R: Yes, but your forms are paradoxical.

Q: How so?

R: Your forms throw out induction. This is simply Hume's problem dressed up, but the solution stays the same. One cannot move forward without some basic axiom of the invariant nature of reality over some configuration space. Primary space and time.

Q: Why can't we move forward?

R: If one questions basic logic, how could one possibly come to any sort of understanding of observations.

Q: As you most certainly know, a lack of an idea is hardly evidence that such an idea is impossible.

R: Ok, ok... If one where to doubt logic, deduction goes out the window. If deduction goes out abduction goes out. If we have no abduction we cannot trust our senses nor our insurements. Thus no evidence, ways to evaluate evidence, nor any theories to understand any evidence. Science goes down with it.

Q: Such a fragile thing science is.

R: I know where you are going. This calls into question the validity of science, hanging on by an unprovable axiom. Yet it matters not. If science works, then it helps, if it doesn't, it has helped in the past and abandoning it would not do us any better.

Q: So pragmatism becomes the answer?

R: In the question of axioms, yes.

Q: Why not everything else?

R: Once we have the basic axioms, such hand waving is unessecary as more rigorous derivations are possible.

Q: Why are these more rigerous methods preferable?

R: We have gone all of the way from spacial invariants to the foundations of epistemology, why don't we give it a rest.

Q: Isn't this what you wanted?

R: What are you going on about, you are me. This meta-humor makes no sense, especially as no one will ever read this.

Q: Do you want someone to read this and think you clever?

R: No, at least I had not thought of it until you(I) wrote that down... I still hold it that it be a personal choice rather than the desire for external validation.

Q: Ok... I won't push it, back to where we where. Why are more rigerous methods perfered in epistemology. If you already concede that absolute truth is unattainable(I know you haven't already outwardly have, but you will get there as you are me and this is quite long[what in the world am I doing, writing to myself?{Yes, that is the whole point((Golly I sound like I had back when I started Logos[[Is that a bad thing, isn't that what you are trying to do?{{Alright, that is enough}}]]))}]) Anyways, wouldn't a pragmatic view be more consistent and make more sense?

R: In many cases, a pragmatic world view can lead to idyllic views where idealism leads to pragmatic outcome. In science this is obviously true.

Q: That is hardly a view in the philosophical sense, but more in the aesthetic.

R: Then I take that on. Pragmatic in epistemology, instramentalism in pragmatic thought, and scientific realist in ideals.

Q: That is quite complicated.

R: It may be, but so is the universe.

\paragraph{Computational Models and their Accuracy*}

\paragraph{Further Development in Structural preserving methods}

\paragraph{Explaining Gravitational Waves}
Q: Explain to me gravitational waves

R: Essentially, gravitational waves, also refereed to as gravitational radiation, are small perturbations in spacetime. Mathematically equivalent to waves.

Q: What causes them?

R: It is something called quadropole moments, where quadropole refers to motion where the distribution of the stress-energy density changes while the ``center of mass'' remains the same. Examples would be binary orbits or CCSN.

Q: Why quadropole?

R: Quadropole moment, refers to distribution, it is a higher dimension extension than monopoles and dipoles. Monopoles referring to changes in the cumulative, such as total mass. While dipoles would be change in center of mass. Obviously, mass is conserved so any change in it would be by its nature a higher dimensional pole. A dipole movement would be uniform and thus unable to preterb spacetime in wave form. While quadropoles would be non-uniform by definition.

Q: What about higher dimensional?

R: I am not sure what higher dimensional poles would be so I would have to look that up.

Q: Before we do that, let us try to examine it without that. Think about an abstraction of what we are looking at.

R: Well, we go from the integral, to the $1/n$ of the integral, to the function itself. The general idea is on the tip of my tounge but I keep think aout it going nowhere.

Q: Ok, look it up, but then tell me about what that unearthed.

R: Multipoles generalize ``shape descriptors.'' Going from a scalar that tells you the total but nothing of movement, symmetry breaking, etc. Then adding movement, then flattening or elongating(quadropoles), then three-fold. They look at the asymmetry of the asymmetry of an earlier order. Taking that all quadpoles plus define the distribution changes while quadropoles alone look at compression and elongating the structure.

Q: Great, what does this say about gravitational waves?

R: Quaropole plus tend to have the heaviest effect while later ones add correction terms.

Q: How do the above examples give these?

R: Take a binary orbit, the distribution of mass obviously changes, the matter appears in different areas. Taking the example of flattening or elongating, it effectively does that, imagine the curvature of spacetime rather than the point particles. This makes the earlier analogy of the sphere much more clear.


\section{Mathematics}

\section{Research Culture}

\section{Broader Culture}


\section{General Conversing}



% LocalWords:  invariants
